% Lösungen Einheit 1 von 2 – Die Volumenformel (zusammenbau v0.1)
\einheitenkopf{Einheit 1 von 2 · Die Volumenformel}

\erg{8}{a) um die y-Achse – der Ansatz passt nicht, es wird über y integriert \quad b) $V = \pi \cdot$ Integral von 0 bis 3 über $(x^2 + 4x + 4)\,dx = 39\pi \approx 122{,}52$ \quad c) $(f(x))^2 = \tfrac{1}{1024}x^4 - \tfrac{3}{16}x^2 + 9$; $V = \pi \cdot$ Integral von −8 bis 8 $= 92{,}8\pi \approx 291{,}5$ cm³ $\approx 0{,}29$ l \quad d) $V_1 = \pi \cdot$ Integral von 0 bis 2 über $(0{,}25x^2 + 2x + 4)\,dx = \tfrac{38}{3}\pi$, $V_2 = \pi \cdot 3^2 \cdot 2 = 18\pi$; zusammen $\tfrac{92}{3}\pi \approx 96{,}34$ dm³, also etwa 96,3 l \quad e) Das Wasser steht mindestens bis zur Kugeloberkante $x = 8$: $\pi \cdot$ Integral von 0 bis 8 über $(4x + 36)\,dx = 416\pi$; Kugel $\tfrac{4}{3}\pi \cdot 4^3 = \tfrac{256}{3}\pi$; Wasser mindestens $416\pi - \tfrac{256}{3}\pi \approx 1\,038{,}8$ cm³ \quad f) $\pi \cdot$ Integral von t bis t + 1 über $(3x + 25)\,dx = 250$ mit $t \ge 6$; die Kugel liegt unterhalb dieser Schicht und wird nicht abgezogen}
\erg{9}{a) Fehler: das gemischte Glied $8x$ fehlt. Richtig: $(f(x))^2 = x^2 + 8x + 16$, $V = \pi \cdot (\tfrac{8}{3} + 16 + 32) = \tfrac{152}{3}\pi \approx 159{,}17$ \quad b) Der Funktionswert an der Stelle x ist der Radius der Kreisscheibe an dieser Stelle, und ihre Fläche ist die Kreiszahl mal Radius im Quadrat. \quad c) $V = \pi \cdot$ Integral von −10 bis 10 über $(16 - 0{,}08x^2 + 0{,}0001x^4)\,dx = \tfrac{812}{3}\pi \approx 850{,}3$ l; 900 Flaschen sind 675 l – ja, es reicht}
